% Setup sketches (2026-10-29), one figure per group of tests that share a geometry.
% Configuration-A channel (L1-1, L1-2, L1-3, L1-3b): existing Figure fig:num_channel_A.
% fig:num_sketch_lattice -> end of subsubsec:num_setup_lattice (L1-4, L1-5, L1-6, L1-8, L1-7)
% fig:num_sketch_B       -> end of subsubsec:num_setup_B (L2-1 ... L2-5; L2-6)
% fig:num_sketch_3D      -> start of subsec:num_3D (L3-1, L3-2)
% Requires the TikZ libraries of main.tex (arrows.meta, calc, decorations.pathmorphing).

%%%% BEGIN fig:num_sketch_lattice
\begin{figure}[t]
    \centering
    \begin{tikzpicture}[x=0.8cm,y=0.8cm,font=\footnotesize,
        spr/.style={decorate,decoration={zigzag,amplitude=1.2pt,segment length=2.5pt}}]
        % (a) finite array in the configuration-A channel
        \node[anchor=west] at (-0.3,1.3) {(a)};
        \fill[blue!6] (0,0) rectangle (16,-1.4);
        \draw[thick] (0,0.2) -- (0,-1.4) -- (16,-1.4) -- (16,0.2);
        \fill[gray!35] (0,0) rectangle (2.5,0.18);
        \fill[gray!35] (13.5,0) rectangle (16,0.18);
        \node[above] at (1.25,0.2) {sponge};
        \node[above] at (14.75,0.2) {sponge};
        \draw[very thick] (0,0) -- (16,0);
        \fill[red!70] (3.9,0) rectangle (4.2,0.15);
        \node[above,red!70!black] at (4.05,0.2) {$F/b$};
        \draw[|-|, blue!70!black] (4.7,0.35) -- (6.3,0.35) node[midway,above] {probe};
        \draw[|-|, blue!70!black] (9.5,0.35) -- (11.1,0.35) node[midway,above] {probe};
        \foreach \x in {6.9,7.4,7.9,8.4,8.9} {
            \draw (\x,0) -- (\x,0.25);
            \draw[spr] (\x,0.25) -- (\x,0.6);
            \fill (\x-0.12,0.6) rectangle (\x+0.12,0.76);
        }
        \node[above] at (7.9,0.8) {$N$ resonators};
        \draw[<->] (6.9,-0.3) -- (7.4,-0.3) node[midway,below] {$a$};
        \draw[dashed] (7.9,0) -- (7.9,-0.75) node[below] {$x_c$};
        \draw[<->] (2.5,-0.5) -- (3.9,-0.5) node[midway,below] {$\ge 5/\operatorname{Im}k_1$};
        \draw[<->] (11.1,-0.5) -- (13.5,-0.5) node[midway,below] {$\ge 5/\operatorname{Im}k_1$};
        \draw[<->] (0.6,0) -- (0.6,-1.4) node[midway,right] {$H$};

        \begin{scope}[shift={(0,-3.7)}]
            % (b) supercell of M cells with Bloch-phased loads
            \node[anchor=west] at (-0.3,1.3) {(b)};
            \fill[blue!6] (0.5,0) rectangle (8.5,-1.4);
            \draw[thick] (0.5,-1.4) -- (8.5,-1.4);
            \draw[thick,dashed] (0.5,0.2) -- (0.5,-1.4);
            \draw[thick,dashed] (8.5,0.2) -- (8.5,-1.4);
            \node[rotate=90,above] at (0.5,-0.7) {periodic};
            \node[rotate=90,below] at (8.5,-0.7) {periodic};
            \draw[very thick] (0.5,0) -- (8.5,0);
            \foreach \i in {1,2,3} \draw[gray] (0.5+2*\i,0.1) -- (0.5+2*\i,-0.1);
            \foreach \x in {1.5,3.5,5.5,7.5} {
                \draw (\x,0) -- (\x,0.25);
                \draw[spr] (\x,0.25) -- (\x,0.6);
                \fill (\x-0.12,0.6) rectangle (\x+0.12,0.76);
                \draw[-{Latex[length=1.5mm]},red!70!black] (\x+0.35,0.65) -- (\x+0.35,0.04);
            }
            \node[above,red!70!black] at (7.85,0.65) {$F\mathrm{e}^{\mathrm{i}qx_n}$};
            \node[above] at (1.5,0.8) {$x_1$};
            \node[above] at (3.5,0.8) {$x_2$};
            \draw[<->] (0.5,-0.35) -- (2.5,-0.35) node[midway,below] {$a$};
            \draw[<->] (0.5,-1.7) -- (8.5,-1.7) node[midway,below] {$Ma$};

            % (c) one-wavelength cell with the homogenised resonator field
            \node[anchor=west] at (9.5,1.3) {(c)};
            \fill[blue!6] (10,0) rectangle (15.6,-1.4);
            \draw[thick] (10,-1.4) -- (15.6,-1.4);
            \draw[thick,dashed] (10,0.2) -- (10,-1.4);
            \draw[thick,dashed] (15.6,0.2) -- (15.6,-1.4);
            \draw[very thick] (10,0) -- (15.6,0);
            \foreach \x in {10.2,10.45,...,15.45} {
                \draw[spr,gray!70!black] (\x,0) -- (\x,0.4);
                \fill (\x-0.06,0.4) rectangle (\x+0.06,0.5);
            }
            \draw[red!70!black,domain=10:15.6,samples=80] plot (\x,{0.85+0.12*sin(deg(2*pi*(\x-10)/5.6))});
            \node[above,red!70!black] at (12.8,0.95) {load $\mathrm{e}^{\mathrm{i}kx}$};
            \node[below] at (12.8,-0.25) {field $q(x)$, density $n_0$};
            \draw[<->] (10,-1.7) -- (15.6,-1.7) node[midway,below] {$L=2\pi/k$};
        \end{scope}
    \end{tikzpicture}
    \caption{Setups for the line lattices (not to scale). (a) Finite array of $N$ resonators, spacing $a$ and centre $x_c$, in the channel of Figure~\ref{fig:num_channel_A} (Section~\ref{subsubsec:num_L1_4}). (b) Closed periodic supercell of $M$ cells with one resonator per cell at $x_n=(n-\tfrac12)a$, each loaded by $F\mathrm{e}^{\mathrm{i}qx_n}$ (Sections~\ref{subsubsec:num_L1_5}--\ref{subsubsec:num_L1_8}); $M=1$ gives the zone centre, $M=2$ with alternating load the zone boundary. (c) Periodic cell of one wavelength, $L=2\pi/k$, carrying the homogenised resonator field $q(x)$ and the load $\mathrm{e}^{\mathrm{i}kx}$ (Section~\ref{subsubsec:num_L1_7}).}
    \label{fig:num_sketch_lattice}
\end{figure}
%%%% END fig:num_sketch_lattice

%%%% BEGIN fig:num_sketch_B
\begin{figure}[t]
    \centering
    \begin{tikzpicture}[x=0.75cm,y=0.8cm,font=\footnotesize,
        spr/.style={decorate,decoration={zigzag,amplitude=1.2pt,segment length=2.5pt}}]
        % (a) free plate in open water
        \node[anchor=west] at (-0.3,1.4) {(a)};
        \fill[blue!6] (0,0) rectangle (16,-1.6);
        \draw[thick] (0,-1.6) -- (16,-1.6);
        \draw[thick,dashed] (0,0.15) -- (0,-1.6);
        \draw[thick,dashed] (16,0.15) -- (16,-1.6);
        \node[rotate=90,above] at (0,-0.8) {Sommerfeld};
        \node[rotate=90,below] at (16,-0.8) {Sommerfeld};
        \draw (0,0) -- (4.5,0);
        \draw (11.5,0) -- (16,0);
        \draw[line width=2.2pt] (4.5,0) -- (11.5,0);
        \draw[-{Latex},thick,blue!70!black] (0.6,0.45) -- (2.6,0.45) node[midway,above] {incident wave};
        \foreach \x/\j in {5.8/1,7.3/2} {
            \draw (\x,0) -- (\x,0.2);
            \draw[spr] (\x,0.2) -- (\x,0.55);
            \draw[-{Latex[length=1.5mm]},red!70!black] (\x,1.05) -- (\x,0.6);
            \node[below] at (\x,-0.05) {$x_\j$};
        }
        \node[above,red!70!black] at (6.55,1.05) {unit force, massless spring};
        \draw (9.8,0) -- (9.8,0.25);
        \draw[spr] (9.8,0.25) -- (9.8,0.6);
        \fill (9.68,0.6) rectangle (9.92,0.76);
        \node[right] at (9.95,0.7) {$m_r,k_r,c_r$};
        \draw[<->] (0,-0.45) -- (4.5,-0.45) node[midway,below] {$\max\{\lambda_\mathrm{out},6/\kappa_1\}$};
        \draw[<->] (11.5,-0.45) -- (16,-0.45) node[midway,below] {$\max\{\lambda_\mathrm{out},6/\kappa_1\}$};
        \draw[<->] (4.5,-1.95) -- (11.5,-1.95) node[midway,below] {$L_p$, free edges};
        \draw[<->] (17.3,0) -- (17.3,-1.6) node[midway,right] {$H$};

        \begin{scope}[shift={(0,-4.6)}]
            % (b) graded array on a long plate with a plate source
            \node[anchor=west] at (-0.3,1.4) {(b)};
            \fill[blue!6] (0,0) rectangle (16,-1.6);
            \draw[thick] (0,-1.6) -- (16,-1.6);
            \draw[thick,dashed] (0,0.15) -- (0,-1.6);
            \draw[thick,dashed] (16,0.15) -- (16,-1.6);
            \draw (0,0) -- (0.6,0);
            \draw (15.4,0) -- (16,0);
            \draw[line width=2.2pt] (0.6,0) -- (15.4,0);
            \fill[gray!35] (0.6,0.08) rectangle (2.8,0.26);
            \fill[gray!35] (13.2,0.08) rectangle (15.4,0.26);
            \node[above] at (1.7,0.26) {sponge};
            \node[above] at (14.3,0.26) {sponge};
            \draw[-{Latex[length=1.5mm]},red!70!black] (3.7,0.9) -- (3.7,0.08);
            \node[above,red!70!black] at (3.7,0.9) {$x_s$};
            \draw[|-|, blue!70!black] (4.3,0.4) -- (5.7,0.4) node[midway,above] {$W_1$};
            \draw[|-|, blue!70!black] (11.2,0.4) -- (12.6,0.4) node[midway,above] {$W_2$};
            \foreach \i in {0,...,14} {
                \pgfmathsetmacro{\x}{6.3+0.3*\i}
                \pgfmathsetmacro{\c}{15+4*\i}
                \draw (\x,0) -- (\x,0.18);
                \draw[spr] (\x,0.18) -- (\x,0.5);
                \fill[black!\c!blue!60] (\x-0.08,0.5) rectangle (\x+0.08,0.62);
            }
            \node[above] at (8.4,0.65) {graded $\omega_r(x)$ or $n(x)$};
            \draw[dashed,orange!80!black] (8.7,0) -- (8.7,-1.0) node[below] {$x^*(\omega)$};
            \draw[|-|,orange!80!black] (8.1,-0.35) -- (9.3,-0.35);
            \node[right,orange!80!black] at (9.3,-0.35) {$\Delta x$};
            \draw[<->] (6.3,-1.95) -- (10.5,-1.95) node[midway,below] {array, $x\in[0,L_g]$};
        \end{scope}
    \end{tikzpicture}
    \caption{Setups for finite structures in open water, configuration~B (not to scale). (a) Free plate of length $L_p$ between open-water regions closed by the Sommerfeld condition. The transfer compliance $G_B$ is obtained from unit forces applied through massless springs at the attachment points $x_j$ (Section~\ref{subsubsec:num_L2_1}). The same domain carries one resonator (Sections~\ref{subsubsec:num_L2_3} and~\ref{subsubsec:num_L2_4}, the latter with an incident wave from the left) or a tuned population with spacing $a\simeq1~\si{\meter}$ (Section~\ref{subsubsec:num_L2_5}). (b) Graded array on a long plate (Section~\ref{subsec:num_graded}): point source $x_s$ on a bare lead-in, fit windows $W_1$ and $W_2$, viscous sponges of $6\lambda_p$ at both plate ends, and all separations at least $5/\kappa_1$. A frequency $\omega$ is absorbed near its turning point $x^*(\omega)$ over a footprint $\Delta x$.}
    \label{fig:num_sketch_B}
\end{figure}
%%%% END fig:num_sketch_B

%%%% BEGIN fig:num_sketch_3D
\begin{figure}[t]
    \centering
    \begin{tikzpicture}[x={(0.8cm,0cm)},y={(0.33cm,0.33cm)},z={(0cm,0.8cm)},font=\footnotesize,
        spr/.style={decorate,decoration={zigzag,amplitude=1.2pt,segment length=2.5pt}}]
        % (a) L3-1: box periodic in y, line load
        \node[anchor=west] at (-1.2,0,1.2) {(a)};
        \draw[gray!60,dashed] (0,3,0) -- (0,3,-2) -- (14,3,-2);
        \draw[gray!60,dashed] (0,3,-2) -- (0,0,-2);
        \fill[blue!6] (0,0,0) -- (14,0,0) -- (14,0,-2) -- (0,0,-2) -- cycle;
        \fill[blue!12] (14,0,0) -- (14,3,0) -- (14,3,-2) -- (14,0,-2) -- cycle;
        \fill[gray!12] (0,0,0) -- (14,0,0) -- (14,3,0) -- (0,3,0) -- cycle;
        \fill[gray!45] (0,0,0) -- (3,0,0) -- (3,3,0) -- (0,3,0) -- cycle;
        \fill[gray!45] (11,0,0) -- (14,0,0) -- (14,3,0) -- (11,3,0) -- cycle;
        \fill[red!60] (5.4,0,0) -- (5.7,0,0) -- (5.7,3,0) -- (5.4,3,0) -- cycle;
        \draw[thick] (0,0,-2) -- (14,0,-2) -- (14,3,-2);
        \draw[thick] (0,0,0) -- (14,0,0) -- (14,3,0) -- (0,3,0) -- cycle;
        \draw[thick] (0,0,0) -- (0,0,-2);
        \draw[thick] (14,0,0) -- (14,0,-2);
        \draw[thick] (14,3,0) -- (14,3,-2);
        \node[above] at (1.5,3,0) {sponge};
        \node[above] at (12.5,3,0) {sponge};
        \node[above,red!70!black] at (5.55,3,0) {$F\mathrm{e}^{\mathrm{i}k_yy}$};
        \draw[-{Latex},thick] (5.7,1.5,0) -- (8.3,2.7,0) node[right] {$\boldsymbol{k}$};
        \draw[-{Latex},thick] (5.4,1.5,0) -- (3.4,2.7,0);
        \node at (7,0,-1) {periodic in $y$ (faces $y=0$ and $y=W$)};
        \draw[<->] (14.4,0,0) -- (14.4,3,0) node[pos=1,right] {$W=\lambda_p/\sin\theta$};
        \draw[<->] (-0.4,0,0) -- (-0.4,0,-2) node[midway,left] {$H$};
        \draw[-{Latex}] (15.6,0,-2.3) -- (16.6,0,-2.3) node[right] {$x$};
        \draw[-{Latex}] (15.6,0,-2.3) -- (15.6,1.2,-2.3) node[right] {$y$};
        \draw[-{Latex}] (15.6,0,-2.3) -- (15.6,0,-1.5) node[above] {$z$};

        \begin{scope}[shift={(0,0,-5.6)}]
            % (b) L3-2: channel with walls, plate and point resonator
            \node[anchor=west] at (-1.2,0,1.6) {(b)};
            \draw[gray!60,dashed] (0,2.4,0) -- (0,2.4,-2) -- (16,2.4,-2);
            \draw[gray!60,dashed] (0,2.4,-2) -- (0,0,-2);
            \fill[blue!6] (0,0,0) -- (16,0,0) -- (16,0,-2) -- (0,0,-2) -- cycle;
            \fill[blue!12] (16,0,0) -- (16,2.4,0) -- (16,2.4,-2) -- (16,0,-2) -- cycle;
            \fill[gray!45] (0,0,0) -- (2.5,0,0) -- (2.5,2.4,0) -- (0,2.4,0) -- cycle;
            \fill[blue!15] (2.5,0,0) -- (4,0,0) -- (4,2.4,0) -- (2.5,2.4,0) -- cycle;
            \fill[gray!15] (4,0,0) -- (12,0,0) -- (12,2.4,0) -- (4,2.4,0) -- cycle;
            \fill[blue!15] (12,0,0) -- (13.5,0,0) -- (13.5,2.4,0) -- (12,2.4,0) -- cycle;
            \fill[gray!45] (13.5,0,0) -- (16,0,0) -- (16,2.4,0) -- (13.5,2.4,0) -- cycle;
            % plate mesh: fixed far field, graded next to the resonator
            \foreach \x in {4.5,5,...,7.5,8.5,9,...,11.5} \draw[gray!85,very thin] (\x,0,0) -- (\x,2.4,0);
            \foreach \x in {7.6,7.75,7.88,8.12,8.25,8.4} \draw[gray!85,very thin] (\x,0,0) -- (\x,2.4,0);
            \foreach \y in {0.4,0.8,1.6,2.0,0.9,1.0,1.1,1.3,1.4,1.5} \draw[gray!85,very thin] (4,\y,0) -- (12,\y,0);
            \draw[thick] (4,0,0) -- (12,0,0) -- (12,2.4,0) -- (4,2.4,0) -- cycle;
            \draw[thick] (0,0,-2) -- (16,0,-2) -- (16,2.4,-2);
            \draw[thick] (0,0,0) -- (16,0,0) -- (16,2.4,0) -- (0,2.4,0) -- cycle;
            \draw[thick] (0,0,0) -- (0,0,-2);
            \draw[thick] (16,0,0) -- (16,0,-2);
            \draw[thick] (16,2.4,0) -- (16,2.4,-2);
            % resonator at the plate centre
            \draw (8,1.2,0) -- (8,1.2,0.3);
            \draw[spr] (8,1.2,0.3) -- (8,1.2,0.9);
            \node[fill,inner sep=2pt] at (8,1.2,0.98) {};
            \node[right] at (8.15,1.2,1.0) {$M_\rho,\omega_r$ at $\boldsymbol{x}_r$};
            \draw[-{Latex},thick,blue!70!black] (-2.2,1.2,0) -- (-0.3,1.2,0) node[pos=0,above right] {incident};
            \node[above] at (1.25,2.4,0) {damping};
            \node[above] at (14.75,2.4,0) {damping};
            \node[above] at (6,2.4,0) {plate, $4\lambda_p$};
            \node at (8,0,-1) {no-flux walls at $y=0$ and $y=W$};
            \draw[<->] (16.4,0,0) -- (16.4,2.4,0) node[pos=1,right] {$W$};
            \draw[<->] (-0.4,0,0) -- (-0.4,0,-2) node[midway,left] {$H$};
        \end{scope}
    \end{tikzpicture}
    \caption{Three-dimensional setups (not to scale). (a) Plate over the whole surface of a box periodic in $y$, with sponges at both ends and a line load $F\mathrm{e}^{\mathrm{i}k_yy}$, $k_y=k_0\sin\theta$, that radiates two oblique plate waves (Section~\ref{subsubsec:num_L3_1}). (b) Channel of width $W$ with no-flux walls: damping zones of $4\lambda_0$, open water of $\lambda_0/2$, and a plate of length $4\lambda_p$ with free edges and a point resonator at its centre $\boldsymbol{x}_r$. The plate cells are fixed far from $\boldsymbol{x}_r$ and graded towards it (Section~\ref{subsubsec:num_L3_2}).}
    \label{fig:num_sketch_3D}
\end{figure}
%%%% END fig:num_sketch_3D
